% % interactnlmsample.tex
% % v1.05 - August 2017
%=================================================================
% \documentclass{interact}
\documentclass[largeformat]{interact}

\usepackage{epstopdf}
\usepackage[numbers,sort&compress]{natbib}

%=================================================================
% PACKAGES
%=================================================================
\usepackage{xurl}           % URL formatting with linebreak support
\usepackage{tabularx}       % Flexible-width tables
\usepackage{float}          % [H] placement specifier
\usepackage{enumitem}       % Customizable list spacing
\usepackage{algorithm}      % Algorithm float environment
\usepackage{algpseudocode}  % Structured pseudocode
\usepackage{orcidlink}      % Clickable ORCID identifiers
\usepackage{subcaption}     % Subfigures and subtables
\usepackage{listings}       % Code typesetting
\usepackage[table]{xcolor}  % Colors + table cell shading (must be single load)
\usepackage{needspace}      % Prevent inappropriate page breaks
\usepackage{multicol}       % Multi-column layout
\usepackage{multirow}       % Multi-row table cells
\usepackage{longtable}      % Multi-page tables
\usepackage{booktabs}       % Professional table rules
\usepackage{array}          % Extended column types
\usepackage{ragged2e}       % Ragged text alignment
\usepackage{calc}           % Width/height calculations
\usepackage{graphicx}       % Image inclusion
\usepackage{tikz}
\usetikzlibrary{arrows.meta, positioning, shapes.multipart, calc}
\usepackage{pgfplots}

\usepgfplotslibrary{groupplots}
\pgfplotsset{compat=1.18}
\usepgfplotslibrary{statistics}

\usepackage{caption}
\captionsetup[figure]{font=footnotesize}
\captionsetup[algorithm]{font=small}

%=================================================================
% LONGTABLE CONFIGURATION
%=================================================================
\setlength{\LTleft}{0pt}
\setlength{\LTright}{0pt}

%=================================================================
% COLOR DEFINITIONS
%=================================================================
\definecolor{rowgray}{gray}{0.96}
\definecolor{ecorrect}{RGB}{226,245,233}
\definecolor{epartial}{RGB}{255,243,205}
\definecolor{ehalluc}{RGB}{248,215,218}
\definecolor{efield}{RGB}{207,226,255}
\definecolor{eparser}{RGB}{233,236,239}
\definecolor{ercp}{RGB}{238,242,255}

%=================================================================
% CUSTOM MACROS AND COLUMN TYPES
%=================================================================
\newcommand{\rcpcell}[1]{\cellcolor{ercp}{#1}}
\newcommand{\errcell}[2]{\cellcolor{#1}#2}
\newcolumntype{T}{>{\ttfamily}l}
\newcolumntype{M}[1]{>{\ttfamily}p{#1}}
\newcolumntype{L}[1]{>{\RaggedRight\arraybackslash}p{#1}}
\newcolumntype{C}[1]{>{\Centering\arraybackslash}p{#1}}

%=================================================================
% BIBLIOGRAPHY CONFIGURATION
%=================================================================
\bibpunct[, ]{[}{]}{,}{n}{,}{,}


%=================================================================
% BEGIN DOCUMENT
%=================================================================
\begin{document}
\articletype{Research Article} 
\title{A Domain-Agnostic Agentic Architecture for Structured Extraction of Engineering Knowledge}
\author{
    \name{
    Oscar Ikechukwu\textsuperscript{a}\thanks{CONTACT Oscar Ikechukwu. Email: oscar.ikechukwu@liu.se},
    Mehdi Tarkian\textsuperscript{a},
    Sanjay Nambiar\textsuperscript{a},
    Marie Jonsson\textsuperscript{a},
    Christoffer Brax\textsuperscript{b}\thanks{Correspondence for the SAAB AB team may be directed to Anton Strömqvist. Email: anton.stromqvist2@saabgroup.com}
    }
    \affil{\textsuperscript{a}Department of Management and Engineering,
    Linköping University, Linköping, Sweden;
    \textsuperscript{b}SAAB AB, Linköping, Sweden}
}
%=================================================================
% \tableofcontents
% \listoftables
% \listoffigures
%=================================================================
\maketitle
%=================================================================
\begin{abstract}
    LLM-based retrieval systems can execute queries and produce fluent answers, yet they rarely preserve a traceable record of what evidence was retrieved, how it was validated, or what reasoning led to synthesis.
    The lack of provenance and auditability limits debugging, compliance verification, and operational trust—critical for engineering applications where correctness depends on interpreting structured specifications within revision-controlled technical documentation. 
    We introduce the Relational Control Plane (RCP), a SQL-backed agentic architecture that persists orchestration state as queryable relational records. RCP enforces a six-stage verify-then-summarise control loop in which synthesis is permitted only after retrieved evidence satisfies validation constraints, transforming potential hallucinations into explicit, auditable failures. Domain behaviour is defined through reusable Strategy and Function libraries, enabling configuration via declarative templates rather than modifications to orchestration logic.
    Evaluation across two industrial domains shows significant improvements: answer correctness increases from 45\% to 85\% with hallucination reduced from 55\% to 5\% in one case study; and from 23\% to 81\% with hallucination reduced from 47\% to 11\% in another, even with a lightweight language model. These results show that reliable agentic AI systems emerge from architectural persistence, structured validation, and traceable workflows rather than model scale.
\end{abstract}
%=================================================================

%=================================================================
\begin{keywords}
Large Language Models; Agentic AI; Multi-Agent Systems; Knowledge Extraction; Retrieval-Augmented Generation; Engineering Informatics; Natural Language Processing 
\end{keywords}
%=================================================================
% \begin{abbreviations}
% AI -- Artificial Intelligence\\
% LLM -- Large Language Model\\
% OCR -- Optical Character Recognition\\
% RAG -- Retrieval-Augmented Generation\\
% RCP -- Relational Control Plane\\
% SQL -- Structured Query Language
% \end{abbreviations}

%==============================================================================
\section{Introduction}
\label{sec:Introduction}
%==============================================================================
% Engineering documentation problem
Finding a single, reliable piece of engineering information rarely involves consulting a single document. Instead, even routine queries require navigating heterogeneous sources:
% Engineering organizations depend on vast, heterogeneous technical documentation including
specifications, drawings, manuals, revision histories, change orders, reports, and standards, dispersed across incompatible systems and repositories. Retrieving precise product, component,
and process details demands cross-referencing sources, interpreting technical notations, tracking supersessions, and reconciling multi-document dependencies—tasks that require tacit domain expertise \cite{hertzum2000information}. As a result, knowledge workers spend roughly 20\% of their time searching for and gathering information \cite{court1997supporting, mckinsey2012social}, often resorting to asking colleagues when systems fail them.
% because navigating heterogeneous document repositories remains prohibitively difficult.

% Engineering changes + LLM promise
Engineering changes consume 30–50\% of product development effort \cite{fricke2000coping}, with individual modifications cascading unpredictably across interconnected components, triggering revisions to drawings, bills of materials, tooling specifications, and test procedures \cite{clarkson2004predicting}. Despite decades of knowledge management research \cite{mcmahon2004knowledge, heisig2010exploring}, the gap between what engineers need and what retrieval systems deliver persists. However, recent advances in large language models (LLMs) have enabled new approaches to extracting and reasoning over such information sources.

% RAG limitations
Retrieval-Augmented Generation (RAG) grounds LLM outputs in retrieved evidence \cite{lewis2020rag}, yet production deployments reveal systematic failure modes
that are particularly acute for engineering modalities such as unit-bearing
specifications, structured tables, and revision-controlled documentation. Hallucination remains structurally inherent to LLM generation \cite{huang2025hallucinationsurvey, kalai2025hallucinate}, and user studies in manufacturing confirm operators still prefer consulting human experts over autonomous LLM-generated answers \cite{kernanfreire2024knowledge}.

% Agentic frameworks insufficient
Agentic frameworks address some of these limitations through iterative reasoning and tool use \cite{yao2023reactsynergizingreasoningacting, wang2024autonomousagents}. However, engineering tasks often require coordination among multiple specialised agents, including table extraction, revision tracking, and standards interpretation. These agents operate over distinct data sources and impose independent validation requirements. Such coordination demands persistent state management across reasoning episodes \cite{packer2023memgpt}, structured validation gates that prevent unsupported synthesis, and fine-grained source attribution linking each claim to specific document locations \cite{rashkin2023attribution}. Existing multi-agent frameworks such as AutoGen \cite{wu2024autogen} and LangGraph provide graph-based orchestration abstractions, but leave persistence, validation, and deterministic replay to application-level implementation.

% Contribution
This paper presents a domain-agnostic agentic framework for multi-agent orchestration in engineering applications. We externalize orchestration state into a SQL-backed relational persistence model, enabling deterministic replay, full observability, and audit-ready traceability. Rather than maintaining transient in-memory execution graphs, all orchestration artifacts—goals, strategies, function calls, parameters, outputs, and validation outcomes—are stored as queryable relational records.

\vspace{1em}
\noindent The contributions of this work are threefold:
\begin{itemize}[label=-]
    \item A persistent Relational Control Plane (RCP) architecture that validates retrieved evidence before synthesis, storing goals, strategies, and validation outcomes as queryable relational records to enable deterministic replay and audit-ready provenance.
    \item A domain-agnostic design based on reusable Strategy and Function libraries, demonstrated across two industrial domains without modifying core orchestration logic.
    \item An empirical evaluation on 200 engineering queries (100 per industrial domain) showing that the RCP framework increases answer correctness from 45\% to 85\% and reduces hallucination from 55\% to 5\% in the hydraulic domain, and from 23\% to 81\% correctness with hallucination reduced from 47\% to 11\% in the aerospace domain, compared to RAG and SQL-backed retrieval baselines.
\end{itemize}

%==============================================================================
\section{Related Work}
\label{sec:related_work}

%==============================================================================
\subsection{Document Understanding for Structured Extraction}
\label{sec:document_understanding}
%==============================================================================
Engineering knowledge is predominantly captured in documents containing heterogeneous elements: dense tables, technical drawings, figures, flowcharts, algorithms, multi-column layouts, and revision annotations. Extracting structured information from these sources requires layout-aware models that preserve spatial and semantic relationships \cite{zhang2025doc}. Surface-level text extraction is insufficient because critical information is also encoded in table structure/notes, cross-cell references, figure captions, measurement units, domain-specific notation etc. \cite{guo2026dtbench, foppiano2019}.


Figure~\ref{fig:doc_type_connections} summarises the heterogeneity of engineering document types encountered in practice, the content modalities they contain, and the extraction challenges each modality poses.
The co-occurrence of multiple modalities within a single document and the non-local dependencies between them (e.g., revision metadata governing table validity, unit annotations anchored to drawing geometry), is a primary reason why surface-level text extraction and generic chunking strategies underperform on engineering knowledge retrieval tasks. See Figure~\ref{fig:sample_table} for an example.

% =====================================================================================
% FIGURE: Engineering document types → content modalities → extraction challenges
% =====================================================================================
\begin{figure}
    \centering
    \includegraphics[width=0.6\linewidth]{figures/document_types.jpg}
    \caption{Heterogeneity of engineering document types and the extraction challenges they pose.}
    \label{fig:doc_type_connections}
\end{figure}


Document AI has advanced through transformer-based architectures and specialised recognition models, including Vision Language Models (VLMs). Xu et al.~\cite{xu2020} introduced LayoutLM for joint text--layout modelling; Table Transformer \cite{lin2022tsrformer, nassar2022tableformer} and  DOTABLER \cite{li2025surface} address table structure recognition; and Safder et al.~\cite{safder2022} demonstrated figure parsing methods for semantic enrichment. However, technical documents introduce additional complexities beyond static layout. Revision markers alter content validity non-locally. A superseded value may appear typographically identical to an active one, requiring applicability metadata to be preserved at extraction time rather than inferred at query time \cite{xie2026agentframework}. Mixed units require normalisation, and cross-references between drawings and parts lists demand anchored, location-aware extraction. Modern multimodal document parsers combine OCR, layout analysis, and entity linking, but often lose semantic anchors such as page numbers, footnotes, references, table cell coordinates, and figure callouts, which are needed for provenance tracking. These challenges motivate extraction pipelines that preserve provenance anchors while transforming unstructured documents into queryable relational structures.


%============================================================================== 
\subsection{Agentic AI in Engineering} 
\label{sec:agentic-ai-engineering} 
%============================================================================== 
Agentic AI systems extend beyond one-shot generation by iteratively planning, invoking external tools, and executing intermediate actions toward a goal \cite{wang2024autonomousagents, prasad2024adapt}. The ReAct paradigm \cite{yao2023reactsynergizingreasoningacting} established interleaved reasoning and acting as a practical template, while multi-agent frameworks such as AutoGen \cite{wu2024autogen} and LangGraph provide graph-based abstractions for coordinating specialised agents. Such systems are increasingly investigated in manufacturing and engineering contexts, where multi-agent coordination and goal-oriented action are required \cite{leng2026, zhao2025}. In engineering document comprehension, agent-based architectures have been proposed to dynamically route queries to specialised agents, motivating modular designs for engineering knowledge extraction \cite{xie2026agentframework}. Empirical work demonstrates both capabilities and limitations: Weibull et al.~\cite{weibull_named_2025} report near-perfect named entity extraction from industrial diagrams using multimodal LLMs, whereas Li et al.~\cite{li2025a} identify persistent issues including specification oversimplification and fabrication when extracting formal specifications from documentation.
 
Two challenges consistently emerge: hallucination and reasoning degradation under extended context. Wei et al.~\cite{wei2026shadows} attribute hallucination to internal-state drift caused by incremental context injection and ``attention-locking'' after multiple interaction rounds. Xie and Ju~\cite{xie2026agentframework} further highlight the limited domain-specific reasoning ability of general-purpose LLMs and argue that constructing curated models for every engineering subdomain is impractical. Packer et al.~\cite{packer2023memgpt} demonstrate the value of externalising memory beyond the context window via a virtual context management layer, though MemGPT does not enforce structured validation gates over retrieved evidence. Madaan et al.~\cite{madaan2023selfrefine} show that iterative self-refinement improves output quality, but their approach operates over free-form language feedback rather than typed, schema-level predicates. Reliability challenges including response inconsistency and limited interpretability were further discussed by Leng et al.~\cite{leng2026}, who conclude that agentic systems require explicit validation mechanisms rather than black-box reasoning. Collectively, these findings motivate architectures that enforce structured interaction, persistent state, and verifiable evidence, allowing intermediate reasoning steps to be inspected and validated, not implicitly assumed.


%==============================================================================
\subsection{Hybrid and Multi-Step Retrieval}
\label{sec:hybrid_multistep_retrieval}
%==============================================================================
RAG \cite{lewis2020rag} is a foundational mechanism for grounding LLM outputs in external knowledge, but has limitations: chunking strategies fragment dense table structures, multi-document consolidation fails when evidence is scattered, and retrieval mechanisms optimised for prose underperform on unit-bearing specifications and structured drawings \cite{barnett2024seven}. These limitations are amplified in engineering
applications requiring synthesis across modalities, document structures, and multiple sources \cite{gao2023retrieval, xie2026agentframework}.

One response has been to pre-structure source knowledge before retrieval. Approaches that objectify unstructured documents into queryable relational or
graph schemas improve retrieval precision by preserving field semantics and enabling constraint-based lookup \cite{aster2024langchainsql}. GraphRAG
\cite{mondal2025} and MDKG-RAG \cite{zhang2025} use graph traversal over structured knowledge graphs; TaxFlow \cite{Karna31122026} combines metadata
indexing with validity filtering for legal texts---analogous to engineering revision control; and HGMem \cite{zhou2026} employs external memory capturing
higher-order correlations to address fragmented reasoning in extended retrieval contexts. Hao et al.~\cite{hao2021} demonstrate that semantically structured
knowledge representations outperform conventional keyword-based search. These approaches strengthen retrieval fidelity but are often architecturally rigid,
since filtering, traversal, and synthesis are largely predefined within the system design.

Parallel work has framed multi-step retrieval as a sequential decision-making problem. Planner--executor architectures (e.g. OPERA-AR \cite{Liu2025OPERAAR}) decompose complex queries into intermediate sub-goals executed iteratively by specialised reasoning and retrieval components. Reinforcement-learning-based approaches \cite{Song2025R1SearcherIT, jin2025searchr, zhang2024beam, sorokin2026qrag} train LLMs to generate intermediate search queries, enabling more adaptive retrieval. While effective, these methods are both computationally expensive and do not provide staged verification at multiple steps, making it difficult to inspect why a chain succeeded, where it failed, or how it should be replayed and audited in engineering settings.

Several systems combine relational storage with LLM reasoning in ways that partially overlap with the proposed approach. LangChain's SQL retriever and LlamaIndex's SQL query engine support direct database querying from natural language \cite{aster2024langchainsql}, but treat SQL as a retrieval mechanism rather than a persistence layer for orchestration state where execution logic remains in transient memory with no multi-level validation. Emerging skill-based abstractions such as declarative Agent Skill specifications \cite{agentskills2024} define modular, discoverable capability manifests, but still do not address execution provenance. The RCP differs from these systems in two respects: (i)~orchestration artefacts (goal definitions, strategy selections, function invocations, and validation outcomes) are first-class persistent records rather than ephemeral execution traces, and (ii)~progression through the reasoning loop is gated by explicit validation predicates at the function, strategy, and goal levels, rather than being left to a single final generation step.

%==============================================================================
\subsection{Evaluating Multi-Step Reasoning}
\label{sec:evaluating_multistep_reasoning}
%==============================================================================
Traditional question--answering benchmarks such as OK-VQA, WebQA, and HotpotQA emphasise final answer accuracy, masking failure modes that are critical for engineering auditability such as retrieval errors and step omission \cite{marino2019okvqa, chang2021webqa, yang2018hotpotqa}. While useful for model comparison, these benchmarks lack the fine-grained, step-level annotations needed to attribute errors to specific retrieval or reasoning stages.

More recent benchmarks address these limitations. Ning et al.~\cite{ning2026mcsearch} introduce MC-SEARCH, a framework for agentic multimodal RAG with step-wise annotated chains spanning five reasoning topologies, proposing HAVE verification with process-level metrics. RAGEval \cite{zhu2025rageval} builds on this direction by introducing Completeness, Hallucination, and Irrelevance metrics for scenario-specific RAG evaluation beyond conventional accuracy measures. However, process metrics alone are insufficient: observability requires persistent execution traces for post-hoc analysis \cite{moshkovich2025agentic}, and controlled field studies show that trust in AI outputs remains a barrier to adoption even when accuracy is high \cite{lowhagen2025aiassistant}. Crucially, none of these benchmarks targets the correctness criteria central to domain-specific agentic retrieval, including unit fidelity, revision applicability, and citation granularity, beyond semantic answer matching \cite{jafari2025measurement}. This motivates the evaluation design in Section~\ref{sec:evaluation}.

%=====================================
\section{Method}
\label{sec:Method}
%=====================================
The proposed framework try to address the requirements identified in Section~\ref{sec:related_work} through a six-stage control loop backed by a relational persistence layer. The following subsections detail the data structures, flow and libraries that support this architecture.
%_________________________________________________________________________________________

\subsection{Storing Engineering Data}
\label{subsec:storing_engineering_data}
To accommodate the heterogeneity of engineering documentation, the framework extracts and organises source material into typed data instances paired with explicit provenance anchors (page, section, paragraph, figure captions, and table-cell references). The resulting storage model follows a two-tier structure: a persistent knowledge base that normalises core engineering data into relational form, and a transient, session-scoped cache used for run-time assembly. As illustrated in Fig.~\ref{fig:harvesting_pipeline}, raw documents are first decomposed into structured content units—such as sections, tables, figures, and associated metadata—before being stored in the \emph{Harvested} database.
 
% Harvested / Structured Database (Objectification Layer)
The \textit{harvested} database acts as the canonical repository for engineering knowledge extracted from source materials. It is stable across sessions and free of reasoning artefacts, supporting incremental extension and reproducible querying. In the domains examined in this work, heterogeneous inputs from PDFs undergo a process termed \emph{objectification}. This ensures that:
\begin{enumerate}[label=(\alph*)]
    \item semantic structure and entities follow a consistent relational schema across documents,
    \item numerical and unit-bearing fields are stored with canonical types,
    \item provenance (page, section, figure, or table anchors) is retained, and
    \item all records are indexed to support deterministic, key-based retrieval.
\end{enumerate}

% ================================================================
% FIGURE: Data harvesting workflow.
% ================================================================
\begin{figure}[H]
    \centering
    \includegraphics[width=0.7\linewidth]{figures/data_extraction.jpg}
    \caption{Data harvesting workflow; Source documents are pre-processed, and data stored in the Harvested database.}
    \label{fig:harvesting_pipeline}
\end{figure}

% Temporary Database (Session Assembly Layer)
During reasoning, strategy plans often require intermediate aggregation, filtering, or transformation of data that should not modify the persistent knowledge base. For this purpose, the framework constructs a temporary, session-scoped database at the start of each reasoning episode (described in Section~\ref{subsec:relational_control_plane}). This store holds:

\begin{enumerate}[label=(\alph*)]
    \item intermediate subsets produced by search or filtering agents,
    \item normalised or reformatted values required for downstream validation, and
    \item transient multi-step artefacts such as assembled evidence sets.
\end{enumerate}

The temporary database is strictly short-lived: its contents exist only for the duration of the reasoning episode and are destroyed at session end. Before destruction, all intermediate artefacts are copied as persistent records into the RCP, ensuring that the complete execution trail is recoverable from the RCP alone.
 
%==============================================================================
\subsection{Relational Control Plane (RCP)}
\label{subsec:relational_control_plane}
%==============================================================================
The RCP serves as the central orchestration database of the proposed framework, externalising coordination logic from code into a normalised relational SQL schema (SQLite for prototyping, PostgreSQL for production). It stores orchestration artefacts (goals, strategies, function invocations, parameters, outputs, and validation outcomes) as first-class relational entities rather than unstructured logs or ephemeral runtime state.
 
The RCP enforces a separation between design-time policy and run-time execution. At design time, \emph{library tables} define the system's capability space: strategies are represented as reusable plan templates, while functions are modelled as atomic operations with declared input--output contracts. Each function's contract is formalised through two companion library tables: \texttt{FunctionParametersLibrary}, which declares the expected input parameters (name, type, and constraints), and \texttt{FunctionOutputLibrary}, which declares the expected output schema (name, type, and postconditions). Together with the \texttt{FunctionLibrary} entry itself, these three tables constitute the complete typed contract for each atomic operation. At run time, corresponding \emph{instance tables} record each concrete goal, strategy, and function invocation, including their parameters, outputs, and validation outcomes. This separation allows behaviour to be authored once and reused across domains while making every orchestration decision a persistent, queryable record.
The RCP relational schema is shown in Fig.~\ref{fig:in-session_instantiation}.


% ================================================================
% FIGURE AND ALGORITHM
% ================================================================
\begin{figure*}[!htbp]
    \centering
    \begin{minipage}[c]{0.35\textwidth}
        \vspace{0pt}
        \centering
        \includegraphics[width=\linewidth]{figures/agentic_flowchart.jpg}
        \caption{Agentic control loop showing six stages: goal, strategy, function, and three validation layers interacting with persistent and temporary databases.}
        \label{fig:agentic_flowchart}
    \end{minipage}%
    \hspace{0.03\textwidth}
    \begin{minipage}[c]{0.6\textwidth}
        \vspace{0pt}
        \begin{algorithm}[H]
            \caption{Agentic Reasoning Loop}
            \label{alg:rcp_loop}
            \small
                \begin{algorithmic}[1]
                \definecolor{linecolor}{RGB}{100,100,100}
                \renewcommand{\alglinenumber}[1]{\color{linecolor}\footnotesize #1:}
                
                \State \textbf{Input:} User Query $Q$, Strategy Library $L_S$, Function Library $L_F$, Parameter Schema Library $L_{FP}$, Output Schema Library $L_{FO}$
                \State \textbf{Output:} Validated Answer $A$ or Failure Reason $E$
                \State \textbf{Validation:} ValidateFunction (per-step), ValidateStrategy (plan-level), ValidateGoal (final acceptance)
                \State
                \State \textbf{Stage 1: Goal definition} \label{alg:stage1}
                \State $G \gets \text{Normalize}(Q)$ \Comment{Determine targets, evidence granularity, \& validation rules}
                \State $\text{RCP.\,LogGoal}(G)$
                \State
                \State \textbf{Stage 2: Strategy selection} \label{alg:stage2}
                \State $S \gets L_S.\text{FindBestMatch}(G)$ \Comment{LLM-guided; see \S\ref{subsec:strategy_selection}}
                \State $\text{RCP.\,LogStrategy}(S)$
                \State
                \State \textbf{Stage 3: Function execution} \label{alg:stage3}
                \State $P \gets S.\text{PlanSteps}$
                \For{each $Step \in P$}
                    \State $f \gets L_F.\text{GetFunction}(Step.\text{Capability})$
                    \State $\sigma_{\text{in}} \gets L_{FP}.\text{GetSchema}(f)$ \Comment{Declared input}
                    \State $\sigma_{\text{out}} \gets L_{FO}.\text{GetSchema}(f)$ \Comment{Declared output}
                    \State \text{Validate}($Step.\text{Inputs},\;\sigma_{\text{in}}$) \Comment{Pre-execution check}
                    \State $O \gets f.\text{Execute}(Step.\text{Inputs})$
                    \State $\text{RCP.\,LogExecution}(f, O)$
                    \State
                    \State \textbf{Stage 4: Function validation} \label{alg:stage4}
                    \If{$\neg \text{ValidateFunction}(O,\;\sigma_{\text{out}})$} \Comment{Post-execution check}
                        \State \textbf{retry} $Step$ with adjusted params.\ \textbf{or break}
                    \EndIf
                \EndFor
                \State
                \State \textbf{Stage 5: Strategy validation} \label{alg:stage5}
                \If{$\neg \text{ValidateStrategy}(S,\;S.\text{FinalOutput})$}
                    \State \textbf{backtrack} to Stage~2 \Comment{Try another strategy}
                \EndIf
                \State
                \State \textbf{Stage 6: Goal validation} \label{alg:stage6}
                \If{$\text{ValidateGoal}(G,\;S.\text{FinalOutput})$}
                    \State \Return $A \gets S.\text{FinalOutput}$
                \Else
                    \State \textbf{backtrack} to Stage~1 \textbf{or} \Return $E$ \Comment{Redefine goal}
                \EndIf
                \end{algorithmic}
            \end{algorithm}
    \end{minipage}
\end{figure*}



Upon receiving a query, the controller instantiates a goal record with explicit success criteria (such as revision scope, evidence granularity, or citation requirements) and selects an appropriate strategy instance according to stored policies in the RCP libraries. The controller then dispatches one or more agents whose input--output contracts match the current system state; all validation outcomes are persisted as structured records in the instance tables.


Progress through the loop is validation-centric: advancement is permitted only when defined predicates such as unit consistency, revision applicability, or citation completeness are satisfied. All computational artefacts (structured tables, retrieved passages, intermediate results, and evidence links) are stored within the RCP, while the orchestration graph maintains only lightweight references (active goal, strategy, and function identifiers). This arrangement minimises in-memory state and enables reproducible runs through standard relational queries. Note that while the artefact trail is fully reproducible, decision steps involving LLM inference (strategy selection, synthesis) may vary across runs due to the inherent non-determinism of language model generation.

Algorithm~\ref{alg:rcp_loop} provides the formal pseudocode description of this method. It formalises three structural properties that distinguish the RCP from conventional agentic frameworks:
\begin{enumerate}
    \item \text{Data-driven policy:} Changes in execution order, eligibility, or capability preference (for example, prioritising hierarchical retrieval over keyword search) are achieved through updates to library records rather than code rewrites.
    \item \text{Full artefact persistence:} Every decision, artefact, and validation outcome is recorded, supporting post-hoc analysis and quantitative evaluation of strategy performance.
    \item \text{Modular extensibility:} New or alternative capabilities (e.g., retrievers, table parsers, or unit normalisation tools) can be registered in the RCP function library; existing strategies incorporate them without altering the orchestration graph.
\end{enumerate}



%==============================================================================
\subsection{Strategy Library}
\label{subsec:strategy_library}
%==============================================================================
The Strategy Library contains reusable plan templates that standardise patterns of tool usage, reasoning, and validation across engineering contexts. Each strategy specifies four core components:
\begin{itemize}[label=- ]
    \item \text{Targets:} The evidence or variables required for successful completion (e.g., numerical value with units, revision identifier, document citation).
    \item \text{Eligible functions:} The permissible computational modules and their preferred execution order (e.g., search $\rightarrow$ parse $\rightarrow$ normalise $\rightarrow$ verify).
    \item \text{Sufficiency conditions:} Formal predicates that must hold before a strategy is considered complete (e.g., at least two corroborating sources; explicit revision supersession confirmed).
    \item \text{Fallback policies:} Adaptive behaviours when sufficiency is not met (e.g., expand search scope, trigger hierarchical retrieval).
\end{itemize}
For example, a \emph{Specification Lookup} strategy declares: target = product attribute with unit; plan = [extract product code $\rightarrow$ retrieve record $\rightarrow$ extract attribute $\rightarrow$ validate]; sufficiency = attribute resolved and unit normalised. Concrete instantiations for each case study are presented in Sections~\ref{subsec:hydroscand_case} and~\ref{subsec:saab_case}.

%==============================================================================
\subsection{Function Library}
\label{subsec:function_library}
%==============================================================================
The Function Library contains modular, atomic operations that execute the steps defined in strategy plans. Each function is a self-contained component with declared inputs, outputs, and postconditions. Functions do not modify the persistent Harvested database; their outputs are written exclusively to session-scoped RCP instance tables, ensuring that the knowledge base remains stable across sessions.
 
Each function's interface is formalised through two companion library tables. \texttt{FunctionParametersLibrary} declares the expected input schema: parameter names, types, and constraints (e.g., \texttt{product\_code: string, required}). \texttt{FunctionOutputLibrary} declares the expected output schema: output names, types, and postconditions (e.g., \texttt{extracted\_data: json, non-empty}). At execution time, the controller validates actual inputs against \texttt{FunctionParametersLibrary} before invocation and validates actual outputs against \texttt{FunctionOutputLibrary} afterwards, enabling early detection of schema violations before downstream functions are invoked.
 
Typical functions include: text retrieval (keyword/semantic search with constraint support), hierarchical retrieval (document tree traversal), table extraction (region detection with cell anchoring), unit normalisation (canonical conversion with provenance), revision checking (supersession tracing), figure linking, keyword expansion, table filtering, and evidence scoring.
 
Each function produces validated artefacts (values, units, revisions, citations) that feed into a strategy's target set. Validation gates determine whether these artefacts satisfy sufficiency conditions; when unmet, the orchestration layer invokes additional or alternative functions. Because all functions adhere to uniform interface specifications, components can be replaced or added by registering them in the Function Library---no changes to the orchestration graph are required.



%==============================================================================
\subsection{In-Session Instantiation}
\label{subsec:in-session_instantiation}
%==============================================================================
At runtime, design-time templates are instantiated as session-scoped records that form the execution ledger. For every user query, the controller creates a \texttt{GoalInSession} record with explicit acceptance criteria. It then selects and materialises a \texttt{StrategyInSession} record linked to the goal via foreign key. Within this strategy context, the controller instantiates one or more functions; each invocation generates three linked record sets:
 
\begin{itemize}[label=-]
    \item \texttt{FunctionInSession}: the invocation record (function identity, strategy link, success status).
    \item \texttt{FunctionParametersInSession}: one row per input parameter, validated against the declared schema in \texttt{FunctionParametersLibrary}.
    \item \texttt{FunctionOutputInSession}: one row per output value, validated against the declared schema in \texttt{FunctionOutputLibrary}.
\end{itemize}
 
\noindent Foreign-key constraints enforce a strict hierarchy: every function instance references a strategy instance, which in turn references a goal instance. This means that any record can be traced back to the originating goal through at most two joins. The orchestration graph itself maintains only lightweight references (active goal, strategy, and function identifiers); all heavy artefacts reside in the RCP instance tables.
 
\begin{figure}[!htbp]
    \centering
    \includegraphics[width=0.7\linewidth]{figures/in-session_instantiation_in_RCP.jpg}
    \caption{In-session instantiation within the RCP. Each user query produces a chain of session-scoped records (\texttt{GoalInSession} $\rightarrow$ \texttt{StrategyInSession} $\rightarrow$ \texttt{FunctionInSession}), with parameter and output records validated against library schemas.}
    \label{fig:in-session_instantiation}
\end{figure}
 
Complete parameter and output records for representative queries are provided in Appendix~\ref{app:detailed_traces}, demonstrating how each function invocation materialises concrete artefacts enabling fine-grained inspection, replay, and failure attribution at record level.

% %==============================================================================
% \subsection{Strategy Selection and System Characteristics}
% \label{subsec:strategy_selection_characteristics}
% %==============================================================================

%==============================================================================
\subsection{Strategy Selection}
\label{subsec:strategy_selection}
%==============================================================================

\subsubsection{Strategy Selection Mechanism}
\label{subsubsec:strategy_selection_mechanism}
The controller selects strategies through an LLM-guided evaluation. Upon instantiating a goal, the controller presents the LLM with the current query, a structured goal description, and the ordered list of available strategy templates (excluding those already attempted in the current session). The LLM is prompted to select the most appropriate strategy based on the following criteria:
\begin{itemize}[label=-]
    \item \text{Entity recognition:} Does the query contain a component or product identifier?
    \item \text{Attribute type:} Is the target a simple lookup (single value), comparison (multiple entities), or aggregation (product family)?
    \item \text{Evidence requirements:} Does the goal require revision awareness, cross-document reasoning, or multimodal synthesis?
    \item \text{Data availability:} Are authoritative documents, sections, or paragraph references identified, or is exploratory search required?
\end{itemize}
 
\noindent The LLM returns a single strategy name drawn from the available list; the selection and its accompanying rationale are persisted to the RCP as a structured record, providing an auditable account of each selection decision. If no valid strategy is returned after retries, the controller falls back to a broad retrieval strategy that combines keyword and semantic search across all available documents, or returns a failure diagnostic. Because selection criteria are embedded in the prompt rather than in hard-coded predicates, new strategies are registered by adding a named template to the Strategy Library, with no changes to the controller.
 
The LLM-guided selection step is inherently non-deterministic: given identical inputs, the selected strategy may vary across runs. This is a known trade-off of the current design. While it provides flexibility and avoids hard-coded routing logic, it means that strategy selection is the one component of the framework that is not fully data-driven. A planned enhancement replaces this step with a rule-based precondition mechanism in which each strategy template declares explicit eligibility predicates, making selection fully deterministic and independently auditable.

\subsubsection{Computational Characteristics}
\label{subsubsec:computational_characteristics}
Under the planned rule-based selection, strategy matching is $O(S \cdot P)$, where $S$ is the number of strategy templates and $P$ is the average number of precondition checks per template; in practice $S < 20$ for domain-specific instantiations, making selection overhead negligible. In the current LLM-guided implementation, selection cost is dominated by a single LLM inference call (typically 2K--4K context tokens).
 
Function execution cost is dominated by external operations: SQL retrieval runs in $O(\log N)$ for indexed lookups to $O(N)$ for unindexed full-text scans over $N$ documents, and LLM inference scales linearly with context length (typically 2K--8K tokens per call with KV-cache). For a session with $G$ goals, $S$ strategy instances, and $F$ function calls, RCP storage is $O(G + S + F \cdot (I + O))$, where $I$ and $O$ are average input and output sizes per function. Measured across both case studies, typical sessions generate 5--20 function calls with aggregate storage under 1~MB per query.


% \begin{figure}[t]
% \centering
% \includegraphics[width=0.85\linewidth]{figures/strategy_selection_flow.png}
% \caption{Illustration of the rule-based strategy selection process. Query-derived goal features are evaluated against strategy template predicates. The first template whose preconditions are satisfied is selected and instantiated for execution.}
% \label{fig:strategy_selection}
% \end{figure}

% \subsubsection{Computational Characteristics}
% \label{subsubsec:computational_characteristics}
% \textbf{Latency characteristics:} End-to-end latency ranges from 2--15 seconds for simple lookups (1--3 function calls) to 30--90 seconds for complex multi-document reasoning (10--20 function calls). Parallel execution (Case~II, Strategy~3) reduces latency by 30--50\% when independent function calls are identified.

% \textbf{Current limitations include:}
%     \item \textbf{Limited cross-session learning:} The framework does not currently adapt strategy selection or function parameters based on historical execution outcomes.

% DOMAIN APPLICATION
%===============================================================
\section{Domain Applications}
\label{sec:domain_applications}
%===============================================================
Two industrial instantiations of the framework are presented below. Case~I (Hydroscand~AB) addresses catalogue-style extraction from dense, multilingual product tables; Case~II (SAAB~AB) targets query-driven reasoning over revision-controlled aerospace documentation. Quantitative evaluation follows in Section~\ref{sec:evaluation}.

% CASE I
%===============================================================
\subsection{Case I: Industrial Hydraulic Products Catalogue (Hydroscand AB)}
\label{subsec:hydroscand_case}
%===============================================================
Industrial hydraulic systems rely on extensive product catalogues that consolidate specifications for hoses, fittings, and coupling assemblies. These documents encode critical attributes, including dimensions, pressure ratings, reinforcement structures, material compositions, and thread standards. The \emph{Hydraulic Produktbok}~\cite{hydroscand2020produktbok}, a widely used Nordic fluid-power reference, is representative: it organises chapter-level product families, construction descriptions, and dense specification tables spanning numerous variants within each product line.

The catalogue's structure poses notable challenges for automated extraction. Variant-level data appear in compact tables with merged cells, multi-row headers, mixed unit conventions, and visually implied hierarchies. Thread standards (BSP/G, JIC, ORFS, NPTF) depend on subtle typographic distinctions that OCR and naive PDF parsers often distort, and free-text descriptions (entirely in Swedish) use specialised terminology that impedes generic retrieval models.

As a result, chunk-based RAG tends to lose table semantics, misinterpret values, or fail to deliver cell-level, engineering-grade citations. Queries based on technical specifications (pressure-rating lookup, thread-compatibility checks, application-based product selection) require domain-grounded retrieval that preserves field structure and provenance.

%===============================================================
\subsubsection{Document Pre-processing and Extraction (Case I)}
\label{subsubsec:pre-processing_case_i}
%===============================================================
Source documents are converted into a normalised relational database (\texttt{harvested.db}). Each page is rendered at high resolution to preserve typographic fidelity, which is critical for distinguishing notation spacing (e.g., ``G~1/2'' vs.\ ``G1/2''), sub/superscripts, and unit annotations that carry semantic meaning.

\begin{figure}[!htbp]
    \centering
    \includegraphics[width=0.9\linewidth]{figures/press_coupling.jpg}
    \caption{Sample data extraction from a product category (press coupling).}
    \label{fig:press_coupling}
\end{figure}

% Objectification for storage
Parsed content is inserted into a relational schema comprising \textit{categories} (e.g., H\"ogtrycksslang, Presskopplingar), \textit{families} (e.g., Kappaflex~2K, T9090), and \textit{product-level variants} (e.g., 1105-10-04, 1059-00-16). The objective is to reconstruct the catalogue's implicit structure: the hierarchical organisation of categories and product families, the geometry of specification tables, and the contextual metadata defining material properties, thread standards, and operating limits.

%_______________________________________________________________________________
\subsubsection{Agentic Reasoning (Case I)}
\label{subsubsec:agentic_reasoning}
The agentic reasoning layer instantiates the six-stage control loop (Section~\ref{subsec:relational_control_plane}) for catalogue-based reasoning. Queries are normalised into goal instances and executed through strategy-selected function plans, with all artefacts persisted in \texttt{agentic.db}.

\noindent \textbf{Strategy templates:}
Product catalogue queries exhibit a small number of recurring patterns. In this case study, most requests fall into four archetypes: (i)~direct specification lookup for a known product code, (ii)~contextual search with attribute constraints, (iii)~compatibility checks between components, and (iv)~aggregation across product families. These archetypes are encoded as declarative strategy templates in the \texttt{StrategyLibrary} table, each specifying a target, an ordered plan skeleton, and a sufficiency condition. Table~\ref{tab:strategy_templates} summarises the strategy templates instantiated for the catalogue domain.

% ==============================================================================
% TABLE: Strategy Templates for Case I
% ==============================================================================
\begin{table*}[!htbp]
\captionsetup{width=0.92\linewidth, skip=2pt}
    \caption{Strategy templates for hydraulic catalogue queries.}
    \label{tab:strategy_templates}
    \centering
    \small
    \setlength{\tabcolsep}{6pt}
    \renewcommand{\arraystretch}{1}
    \begin{tabularx}{\textwidth}{c p{1.9cm} X p{4.5cm} p{2.8cm}}
    \toprule
    \textbf{ID} &
    \textbf{Strategy} &
    \textbf{Description} &
    \textbf{Plan Steps} &
    \textbf{Suff.\ Condition} \\
    \midrule
    1 & 
    \textit{Specification Lookup} &
    Retrieve a known specification from an identified product code. &
    Extract product code $\rightarrow$ Query database $\rightarrow$ Extract attributes $\rightarrow$ Analyse with LLM &
    Attribute found and unit normalised. \\
    \midrule
    2 &
    \textit{Contextual Product Search} &
    Identify candidate products matching textual or semantic constraints. &
    Extract requirements $\rightarrow$ Semantic search $\rightarrow$ Filter items $\rightarrow$ Extract attributes $\rightarrow$ Analyse with LLM &
    At least one valid product returned. \\
    \midrule
    3 &
    \textit{Compatibility Check} &
    Assess coupling or thread compatibility between two hydraulic components. &
    Extract component specifications $\rightarrow$ Match standards $\rightarrow$ Check dimensional compatibility $\rightarrow$ Analyse with LLM &
    Standards matched; no dimensional conflict. \\
    \midrule
    4 &
    \textit{Family Summary} &
    Aggregate specification data across all variants within a product family. &
    Retrieve family variants $\rightarrow$ Extract attributes $\rightarrow$ Aggregate results $\rightarrow$ Analyse with LLM &
    Full variant set retrieved and schema consistent. \\
    \bottomrule
    \end{tabularx}
\end{table*}
% ==============================================================================

\noindent \textbf{Function library:} Table~\ref{tab:agent_functions_hydraulic_catalogue} lists sample functions used in this domain, each adhering to a standardised interface (\texttt{success:~bool, result:~dict}).

% ==============================================================================
% TABLE: Function library for the hydraulic catalogue domain.
% ==============================================================================
\begin{table*}[!htbp]
\captionsetup{width=0.92\linewidth, skip=2pt}
    \caption{Function library for the hydraulic catalogue domain.}
    \label{tab:agent_functions_hydraulic_catalogue}
    \centering
    \small
    \setlength{\tabcolsep}{5pt}
    \renewcommand{\arraystretch}{1}
    \begin{tabularx}{\textwidth}{c >{\raggedright\arraybackslash}p{4cm} p{5.8cm} X}
    \toprule
    \textbf{ID} & 
    \textbf{Function Name} & 
    \textbf{Purpose} &
    \textbf{Artefact Produced} \\
    \midrule
    1 & \textit{func\_extract\_products} &
    Retrieve candidate products using keywords and filters. &
    List of product codes with summary specifications. \\
    \midrule
    2 & \textit{func\_filter\_by\_attribute} &
    Narrow a product list by applying attribute constraints (e.g., pressure, DN size). &
    Filtered list satisfying the specified constraints. \\
    \midrule
    3 & \textit{func\_get\_product\_details} &
    Retrieve complete specifications and provenance for a given product. &
    JSON object with metadata and catalogue page reference. \\
    \midrule
    4 & \textit{func\_extract\_attributes} &
    Extract target attribute values from a product specification. &
    Attribute values with units and citations. \\
    \midrule
    5 & \textit{func\_normalize\_units} &
    Convert values to canonical units while preserving original formats. &
    \(420\,\text{bar} \rightarrow 42\,\text{MPa}\). \\
    \midrule
    6 & \textit{func\_analyse\_with\_llm} &
    Compose the final answer with citations and verified source trace. &
    Grounded explanation linked to catalogue data. \\
    \midrule
    7 & \textit{func\_extract\_product\_code} &
    Identify and validate a product identifier from a free-text query. &
    Validated product code (e.g., 4201-16-16). \\
    \bottomrule
    \end{tabularx}
\end{table*}
% ==============================================================================

\textit{func\_analyse\_with\_llm} performs structured LLM synthesis over verified intermediate artefacts. It is instructed to preserve catalogue terminology (e.g., ``Arbetstryck'', ``G\"angtyp''), avoid unsupported claims, and reference source-provenance fields. The remaining functions operate deterministically over structured records in \texttt{harvested.db}, enabling validation gates to intercept failures before synthesis.

% ==============================================================================
\subsubsection{Worked Example: End-to-End Execution (Case I)}
\label{subsubsec:worked_example_case_I}
To illustrate end-to-end behaviour, this subsection traces a representative specification query:
\begin{quote}
    % \emph{Vilken g\"angstorlek har en 4201-16-16?} 
    \emph{``What is the thread size of 4201-16-16?''}
\end{quote}

\noindent To answer this query, the system must: \textbf{(i)}~recognise the product code, \textbf{(ii)}~retrieve the corresponding catalogue record, \textbf{(iii)}~extract the requested attribute (thread size), and \textbf{(iv)}~generate a grounded response.
\vspace{1em}

% ==============================================================================
\noindent \textbf{Goal definition:}
Execution begins at \texttt{GoalDefine}, where the free-text query is normalised into a structured goal instance with explicit evidence requirements. An LLM extracts key entities (product: ``4201-16-16''; attribute: ``g\"angstorlek'') and the controller writes a new record to \texttt{GoalInSession}. The goal specifies the target type (product specification) and a validation predicate requiring that the thread-size attribute is resolved for the identified product before the session can be marked successful (Table~\ref{tab:worked_example_goal}).

% ==============================================================================
\noindent \textbf{Strategy selection:}
The goal instance is matched against \texttt{StrategyLibrary}. Because the query contains an explicit high-confidence product identifier, the controller selects \textit{Specification Lookup} rather than broader retrieval strategies. The resulting strategy instance, including an ordered plan (\texttt{PlanSteps}) and sufficiency condition, is persisted to \texttt{StrategyInSession} (Table~\ref{tab:strategy_instance_table}).

% ==============================================================================
% TABLE: Goal and strategy instances for the worked example (Case I).
% ==============================================================================
\begin{table}[!htbp]
    \captionsetup{width=0.95\linewidth, skip=5pt}
    \caption{Goal and strategy instances for the worked example (Case I).}
    \label{tab:goal_strategy_sidebyside}
    \centering
    \small
    \setlength{\tabcolsep}{2pt}
    \renewcommand{\arraystretch}{1.0}
    \subcaptionsetup{skip=0.5pt}
    \subcaptionbox{Goal instance.\label{tab:worked_example_goal}}[0.49\linewidth]{
        \begin{tabular}{@{}p{0.32\linewidth}@{\hspace{2pt}}p{0.64\linewidth}@{}}
            \toprule
            \textbf{Field} & \textbf{Value} \\
            \midrule
            GoalID & \texttt{1} \\
            SessionID & \texttt{083} \\
            GoalName & \texttt{Main goal} \\
            Target & \texttt{Product specification} \\
            Validation & \texttt{Thread size resolved for identified product} \\
            Description & \texttt{Determine the standard thread size of a 4201-16-16 coupling} \\
            GoalSuccess & \texttt{NULL} \\
            \bottomrule
        \end{tabular}
        }
        \hfill
        \subcaptionbox{Strategy instance.\label{tab:strategy_instance_table}}[0.49\linewidth]{
        \begin{tabular}{@{}p{0.34\linewidth}@{\hspace{3pt}}p{0.60\linewidth}@{}}
            \toprule
            \textbf{Field} & \textbf{Value} \\
            \midrule
            StrategyID & \texttt{1} \\
            GoalID & \texttt{1} \\
            StrategyName & \texttt{Specification Lookup} \\
            Target & \texttt{Product specification} \\
            Description & \texttt{Direct lookup of a known specification from an identified product code} \\
            PlanSteps & \texttt{Extract product code; Retrieve product details; Extract attributes; Analyse with LLM} \\
            Sufficiency & \texttt{Attribute found and unit normalised} \\
            StrategySuccess & \texttt{NULL} \\
            \bottomrule
        \end{tabular}
    }
\end{table}

% ==============================================================================
\noindent \textbf{Function execution and validation:}
Execution proceeds through four function calls (Tables~\ref{tab:worked_example_func1}--\ref{tab:worked_example_func4}): product code extraction, database lookup, attribute extraction, and LLM synthesis. After each step, \texttt{FunctionValidate} enforces basic correctness—non-empty output, expected keys present, count~$>$~0—before the plan advances. Failure at any stage routes control back to \texttt{StrategyValidate} rather than allowing the LLM to synthesize an answer from incomplete evidence.

%==========================================================
% TABLE: Function execution trace..
%==========================================================
\begin{table}[!htbp]
    \captionsetup{width=0.8\linewidth, skip=5pt}
    \caption{Function execution trace.}
    \label{tab:worked_example_grouped}
    \centering
    \small
    \setlength{\tabcolsep}{2pt}
    \renewcommand{\arraystretch}{1.0}
    
    \subcaptionsetup{skip=0.5pt}
    \subcaptionbox{Extract Product Number.\label{tab:worked_example_func1}}[0.49\linewidth]{
    \begin{tabular}{@{}p{0.34\linewidth}@{\hspace{3pt}}p{0.60\linewidth}@{}}
        \toprule \textbf{Field} & \textbf{Value} \\
        \midrule
        FunctionID & \texttt{1} \\
        StrategyID & \texttt{1} \\
        FunctionName & \texttt{func\_extract\_products} \\
        Inputs & \texttt{Query string} \\
        Outputs & \texttt{Keyword Output = 4201-16-16} \\
        ValidationResult & \texttt{Output non-empty; format valid} \\
        FunctionSuccess & \texttt{1} \\
        \bottomrule
    \end{tabular}
    }\hfill
    \subcaptionbox{Query Database.\label{tab:worked_example_func2}}[0.49\linewidth]{
    \begin{tabular}{@{}p{0.34\linewidth}@{\hspace{3pt}}p{0.60\linewidth}@{}}
        \toprule \textbf{Field} & \textbf{Value} \\
        \midrule
        FunctionID & \texttt{2} \\
        StrategyID & \texttt{1} \\
        FunctionName & \texttt{func\_get\_product\_details} \\
        Inputs & \texttt{Keyword Output = 4201-16-16} \\
        Outputs & \texttt{Product record; count = 1} \\
        ValidationResult & \texttt{count $>$ 0; count and format valid} \\
        FunctionSuccess & \texttt{1} \\
        \bottomrule
    \end{tabular}
    }
    
    \vspace{15pt}
    
    \subcaptionbox{Extract Attributes.\label{tab:worked_example_func3}}[0.49\linewidth]{
    \begin{tabular}{@{}p{0.34\linewidth}@{\hspace{3pt}}p{0.60\linewidth}@{}}
        \toprule \textbf{Field} & \textbf{Value} \\
        \midrule
        FunctionID & \texttt{3} \\
        StrategyID & \texttt{1} \\
        FunctionName & \texttt{func\_extract\_attributes} \\
        Inputs & \texttt{Product specification record} \\
        Outputs & \texttt{Extracted attributes: \textit{gängstorlek} = G1"} \\
        ValidationResult & \texttt{Attribute present and non-empty} \\
        FunctionSuccess & \texttt{1} \\
        \bottomrule
    \end{tabular}
    }\hfill
    \subcaptionbox{Analyse with LLM.\label{tab:worked_example_func4}}[0.49\linewidth]{
    \begin{tabular}{@{}p{0.34\linewidth}@{\hspace{3pt}}p{0.60\linewidth}@{}}
        \toprule \textbf{Field} & \textbf{Value} \\
        \midrule
        FunctionID & \texttt{4} \\
        StrategyID & \texttt{1} \\
        FunctionName & \texttt{func\_analyse\_with\_llm} \\
        Inputs & \texttt{Query and extracted attributes} \\
        Outputs & \texttt{Natural-language answer stating thread size} \\
        ValidationResult & \texttt{Output non-empty; well-formed} \\
        FunctionSuccess & \texttt{1} \\
        \bottomrule
    \end{tabular}
    }
\end{table}
%==========================================================

\noindent \textbf{Function parameter/output schemas:}
\texttt{FunctionParametersLibrary} and \texttt{FunctionOutputLibrary} define expected keys, types, and semantics for each registered function. At execution time, these schemas validate records written to \texttt{FunctionParametersInSession} and \texttt{FunctionOutputInSession}, enabling early detection of malformed inputs or missing fields before downstream functions are invoked.

Each function invocation results in multiple parameter and output records, materialized as rows in the in-session instantiation tables. Complete records for this worked example are provided in Appendix~\ref{app:detailed_traces_case1}.


% CASE II
%===============================================================================
\subsection{Case II: Revision-Controlled Aerospace Documentation (SAAB AB)}
\label{subsec:saab_case}
%===============================================================================
The second case study addresses revision-controlled aerospace documentation in collaboration with SAAB~AB (Svenska Aeroplan Aktiebolaget). Unlike Case~I, where queries target a well-defined product catalogue, Case~II requires cross-document reasoning with applicability constraints, revision scope resolution, and multimodal evidence fusion. A single query may require coordinated retrieval from textual, tabular, and graphical elements across several documents.

In compliance with confidentiality requirements, proprietary identifiers, document titles, part numbers, and content values have been anonymised or abstracted. The execution traces, schemas, and control logic presented here nevertheless reflect real runtime artefacts persisted within the RCP.

%================================================================
\subsubsection{Document Pre-processing and Extraction (Case II)}
\label{subsubsec:saab_document_harvesting}
%================================================================
The extraction pipeline produces a document-centric knowledge database that preserves applicability metadata and structural context. The ingested corpus comprises 54 source documents yielding 451 technical tables, 576 annotated images, and 3{,}488 text segments. The domain encompasses six product series of component variants spanning electronic connectors, cable assemblies, and filter components, with product codes following multiple format conventions requiring dedicated normalisation (e.g., space/dash insertion, prefix handling). During ingestion, each document is represented as linked relational units capturing: (i)~document identity and revision markers, (ii)~section hierarchy, and (iii)~modality-specific artefacts (text passages, table fragments, and figure-referenced elements).


\begin{figure}
    \centering
    \includegraphics[width=0.7\linewidth]{figures/sample_tabular_document.jpg}
   \caption{Sample tabular specification data requiring structured extraction.}
    \label{fig:sample_table}
\end{figure}




% \begin{figure}[t]
% \centering

% \begin{subfigure}{0.6\linewidth}
%     \centering
%     \includegraphics[width=\linewidth]{figures/document_types.jpg}
%     \caption{Heterogeneous engineering document types.}
%     \label{fig:doc_type_connections}
% \end{subfigure}

% \vspace{0.5em}

% \begin{subfigure}{0.6\linewidth}
%     \centering
%     \includegraphics[width=\linewidth]{figures/sample_tabular_document.jpg}
%     \caption{Sample engineering tabular data.}
%     \label{fig:sample_table}
% \end{subfigure}

% \caption{Examples of engineering document heterogeneity and structured tabular content. (a) Diverse document types encountered in engineering workflows. (b) Example of tabular specification data requiring structured extraction.}
% \label{fig:doc_combined}

% \end{figure}



Extracted records retain provenance anchors (section labels, table captions, page references), keeping evidence selection location-aware and supporting disambiguation when overlapping revisions coexist. Sample table needing structuring is shown in Figure~\ref{fig:sample_table}.

%================================================================
\subsubsection{Agentic Reasoning (Case II)}
\label{subsubsec:saab_agentic_reasoning}
%================================================================
Case~II reuses the same six-stage control loop without architectural modification, instantiated through a richer strategy and function set tailored to document-centric reasoning. Compared to Case~I, strategies here involve deeper plan structures (up to six sequential steps), parallel execution blocks, and multimodal artefacts including figure-referenced layout generation.

%================================================================
\noindent \textbf{Strategy templates:}
Table~\ref{tab:saab_strategy_templates} presents strategy templates for this domain. Each template follows the structure defined in Section~\ref{subsec:strategy_library}: targets, eligible functions, sufficiency conditions, and fallback policies.

%==========================================================
% TABLE: Strategy templates instantiated for case II domain.
%==========================================================
\begin{table*}[!htbp]
\captionsetup{width=0.92\linewidth, skip=2pt}
    \caption{Strategy templates instantiated for Case II domain.}
    \label{tab:saab_strategy_templates}
    \centering
    \small
    \setlength{\tabcolsep}{6pt}
    \renewcommand{\arraystretch}{1.1}
    \begin{tabularx}{\textwidth}{c p{2.2cm} X p{5cm} X}
    % \begin{tabularx}{\textwidth}{>{\centering\arraybackslash}p{0.5cm} p{1.5cm} X p{4.2cm} p{3.5cm}}
    
    \toprule
    \textbf{ID} &
    \textbf{Strategy} &
    \textbf{Description} &
    \textbf{PlanSteps} &
    \textbf{Suff. condition} \\
    
    \midrule
    
    1 & \textit{Simple Lookup} &
    Deterministic retrieval of structured information from a single authoritative document. & Extract product code $\rightarrow$ Find latest document $\rightarrow$ Table search
    $\rightarrow$ Filter table $\rightarrow$ Assemble table $\rightarrow$ Analyse with LLM & Requested attribute resolved with supporting evidence. \\
    
    \midrule
    
    2 & \textit{Enhanced Lookup} &
    Multi-stage table reasoning with contextual filtering and aggregation. &
    Extract product code $\rightarrow$ Table search $\rightarrow$ Filter table
    $\rightarrow$ Assemble table $\rightarrow$ Analyse with LLM & Consistent result derived across candidate records. \\
    
    \midrule
    
    3 & \textit{Parallel Enhanced Lookup} &
    Latency-aware strategy executing independent steps concurrently. &
    Extract product code $\rightarrow$ Find latest document $\rightarrow$
    [Table search $\parallel$ Document search]
    $\rightarrow$ Filter table $\rightarrow$
    [Analyse with LLM $\parallel$ Convert units]
    $\rightarrow$ Aggregate results & Requested attribute resolved and validated. \\
    
    \midrule
    
    4 & \textit{Visual Layout Generation} &
    Multimodal retrieval and synthesis of layout-level information. &
    Extract product code $\rightarrow$ Retrieve image asset
    $\rightarrow$ Generate visual layout $\rightarrow$ Analyse with LLM &
    Layout produced and grounded in source document. \\
    
    \bottomrule
    \end{tabularx}
    
    \vspace{2mm}
    \raggedright
    \footnotesize
    \textit{Note:} Parallel steps are indicated with the $\parallel$ symbol.
\end{table*}
%================================================================

\noindent \textbf{Function library:}
Table~\ref{tab:saab_function_library} lists some sample functions used in this domain.

% ==========================================================
% TABLE: Functions used in the SAAB documentation case.
% ==========================================================
\begin{table*}[!htbp]
\captionsetup{width=0.92\linewidth, skip=2pt}
    \caption{Functions used in the Case II.}
    \label{tab:saab_function_library}
    \centering
    \small
    \setlength{\tabcolsep}{5pt}
    \renewcommand{\arraystretch}{1.1}
    \begin{tabularx}{\textwidth}{c >{\raggedright\arraybackslash}p{4.2cm} p{6.0cm} X}
    \toprule
    \textbf{ID} &
    \textbf{Function Name} &
    \textbf{Purpose} &
    \textbf{Artifact Produced} \\
    \midrule
    1 & \textit{func\_extract\_product\_code} &
    Identify and validate a product or component identifier from a free-text query. &
    Validated identifier token. \\
    \midrule
    2 & \textit{func\_find\_latest\_document} &
    Select the most recent applicable document revision from the corpus. &
    Document surrogate ID. \\
    \midrule
    3 & \textit{func\_table\_search} &
    Locate candidate tables within a document matching query constraints. &
    Table records and metadata. \\
    \midrule
    4 & \textit{func\_filter\_table} &
    Apply rule-based and semantic filters to tabular data. &
    Filtered table rows. \\
    \midrule
    5 & \textit{func\_assemble\_table} &
    Merge and normalize table fragments into a unified structure. &
    Structured table artifact (JSON). \\
    \midrule
    6 & \textit{func\_generate\_visual\_layout }&
    Produce a schematic or layout-level representation from document figures. &
    Rendered layout artifact. \\
    \midrule
    7 & \textit{func\_analyse\_with\_llm}&
    Synthesize intermediate artifacts into a grounded natural-language response. &
    Final answer text with trace context. \\
    \bottomrule
    \end{tabularx}
\end{table*}
%==========================================================

\subsubsection{Worked Example: End-to-End Execution (Case II)}
We trace a representative query through the \emph{Simple Lookup} strategy. The query targets a specific engineering attribute for an electronic connector component; proprietary identifiers are masked.

Execution proceeds through six function calls: product code extraction, document-revision resolution (selecting the most recent applicable revision from five candidate documents), table search (scanning 451 tables and returning two candidates), table filtering (narrowing to a single matching row), table assembly, and LLM-based synthesis.

%==========================================================
\noindent \textbf{Goal definition:}
A free-text query is normalised into a structured goal instance at the \texttt{GoalDefine} node. The goal specifies the target attribute and a validation predicate requiring that the attribute be resolved with evidence from an authoritative document.

%==========================================================
\noindent \textbf{Strategy selection:}
The controller selects the \emph{Simple Lookup} strategy based on the presence of a high-confidence component identifier. The resulting strategy instance binds the template plan to the current goal.
%==========================================================

\begin{table}[!htbp]
    \captionsetup{width=0.95\linewidth, skip=5pt}
    \caption{Goal and strategy instances (Case II).}
    \label{tab:saab_goal_strategy}
    \centering
    \small
    \setlength{\tabcolsep}{2pt}
    \renewcommand{\arraystretch}{1.0}
    \subcaptionsetup{skip=0.5pt}
    \subcaptionbox{Goal instance.\label{tab:saab_goal_instance}}[0.49\linewidth]{
        \begin{tabular}{@{}p{0.32\linewidth}@{\hspace{2pt}}p{0.64\linewidth}@{}}
            \toprule
            \textbf{Field} & \textbf{Value} \\
            \midrule
            GoalID & \texttt{G\_01} \\
            SessionID & \texttt{S\_XX} \\
            Target & \texttt{Engineering attribute} \\
            Validation & \texttt{Attribute resolved with document provenance} \\
            Description & \texttt{Determine attribute X for component COMP\_[ID-01]} \\
            GoalSuccess & \texttt{NULL} \\
            \bottomrule
        \end{tabular}
        }
        \hfill
        \subcaptionbox{Strategy instance.\label{tab:saab_strategy_instance}}[0.49\linewidth]{
        \begin{tabular}{@{}p{0.32\linewidth}@{\hspace{2pt}}p{0.64\linewidth}@{}}
            \toprule
            \textbf{Field} & \textbf{Value} \\
            \midrule
            StrategyID & \texttt{STR\_01} \\
            GoalID & \texttt{G\_01} \\
            StrategyName & \texttt{Simple Lookup} \\
            PlanSteps & \texttt{Extract product code; Find latest document; Table search; Filter table; Assemble table; Analyse with LLM} \\
            Sufficiency & \texttt{Attribute resolved with supporting evidence} \\
            StrategySuccess & \texttt{NULL} \\
            \bottomrule
        \end{tabular}
    }
\end{table}
%==========================================================

\noindent \textbf{Function execution and validation:}
Execution proceeds through the ordered function plan. After each invocation, \texttt{FunctionValidate} enforces schema-level correctness and expected artefact presence before allowing progression.

%==========================================================
\noindent \textbf{Function parameter/output schemas:}
Sample execution traces with parameter and output schemas are provided in Appendix~\ref{app:detailed_traces_case2}. Tables~\ref{tab:saab_output_excerpt} and~\ref{tab:saab_parameter_excerpt} present abridged excerpts showing how function invocations are recorded in \texttt{FunctionOutputInSession} and \texttt{FunctionParametersInSession}.



% EVALUATION
%============================================
\section{Evaluation}
\label{sec:evaluation}
%============================================
\subsection{Experimental Setup}
\label{subsec:experimental_setup}
%============================================
We evaluate the proposed RCP-based agentic framework on two industrial case studies using an annotated query set of 100 technical questions per case study (200 queries total). Queries were authored by domain engineers at the respective companies who had no involvement in framework development; ensuring separation between query construction and system design. To isolate the contribution of each architectural layer, we compare three systems in a layered comparison that incrementally adds architectural components:

\begin{enumerate}[label=\textit{B\arabic*:}, leftmargin=3.5em]
    \item \textit{Naive RAG.} Standard retrieval-augmented generation: chunk source text (800~characters, 100-character overlap), embed with \textit{qwen3-embedding:8b}\footnote{Qwen3 embedding model, publicly available via Ollama: \url{https://ollama.com/library/qwen3-embedding}.}, store in an in-memory ChromaDB vectorstore, retrieve top-5 chunks, and generate with a single LLM call. No structured schema, no multi-step reasoning.
    \item \textit{SQL-backed Retrieval.} Single-step retrieval over the structured \texttt{harvested.db} schema (categories $\rightarrow$ families $\rightarrow$ products). Product codes are extracted from the query via pattern matching, and direct SQL joins retrieve complete specification records with provenance (page, family, category). Results are passed to the same LLM in a single generation step. This baseline isolates the value of \emph{objectification} (structured relational storage) without multi-step reasoning or validation.
    \item \textit{RCP Framework (Proposed).} The full six-stage control loop with persistent relational state, strategy selection, iterative function execution, and multi-level validation gates.
\end{enumerate}

All methods operate over the same harvested document corpus (\texttt{harvested.db}) and use \emph{Llama3.2:8b} (Ollama~v0.17.7, temperature~$= 0$) for all generation tasks. The compact 8B-parameter model was chosen to reflect resource-constrained industrial deployment conditions where cloud-hosted commercial APIs are precluded by data confidentiality requirements. The Qwen3 embedding model was selected for B1 as a competitive open-weight encoder available through the same local inference stack.

For each case study, we report corpus scale (pages/tables), evaluation scale (number of queries), and per-query compute footprint (function calls and end-to-end runtime). Per-query token counts were not instrumented in the current implementation; token-level cost profiling is planned for future work. Implementation code is available via the Data Availability Statement.

%============================================
\subsection{Quantitative Metrics}
\label{subsec:quantitative_metrics}
%============================================
We report five metrics across both case studies:
(a)~\textit{answer correctness}:binary judgement against the annotated reference set, where a response is scored as correct only if it resolves the target attribute with the exact value and appropriate qualification;
(b)~\textit{citation accuracy}: whether the cited source location (passage, table cell, or page reference) contains the claimed value;
(c)~\textit{unit fidelity}: whether numerical values carry correct units, including normalisation and conversions where applicable;
(d)~\textit{hallucination rate}: the fraction of responses containing at least one factual claim unsupported by retrieved evidence; and
(e)~\textit{operational cost}, measured in end-to-end latency per query.

%============================================
\subsubsection{Performance Results}
\label{subsubsec:performance_results}
%============================================
The Case~I corpus spans approximately 50 catalogue pages encoding 1{,}628 product variants across 168 families; the 100-query set covers direct specification lookups, multi-product comparisons, and contextual product searches. The Case~II corpus comprises 451 extracted tables from 54 revision-controlled aerospace PDFs; the 100-query set covers direct specification lookups and multi-product comparisons across connector families (RPT, RNT, TFR, C0-series). Table~\ref{tab:performance_combined} presents results for all three systems across both domains.

%================================================
% TABLE: Combined quantitative performance comparison
%================================================
\begin{table}[h]
    \centering
    \caption{Quantitative performance comparison across both case studies ($n = 100$ queries each). B1: Naive RAG; B2: SQL-backed Retrieval; B3: RCP (Proposed). \textbf{Bold} = best per row; $\downarrow$ = lower is better. All values in \%.}
    \label{tab:performance_combined}
    \small
    \setlength{\tabcolsep}{5pt}
    \renewcommand{\arraystretch}{1.1}
    \begin{tabular}{@{}lrrrrrr@{}}
        \toprule
        & \multicolumn{3}{c}{\textbf{Case~I}} & \multicolumn{3}{c}{\textbf{Case~II}} \\
        \cmidrule(lr){2-4}\cmidrule(lr){5-7}
        \textbf{Metric} & \textbf{B1} & \textbf{B2} & \textbf{B3} & \textbf{B1} & \textbf{B2} & \textbf{B3} \\
        \midrule
        Answer Correctness (\%)              & 45.0 & 73.0 & \textbf{85.0} & 23.0 & 76.0 & \textbf{81.0} \\
        Citation Accuracy (\%)              & \textbf{99.0} & 93.0 & 90.0 & 33.0 & 59.0 & \textbf{78.0} \\
        Unit Fidelity (\%)                  & 67.0 & \textbf{93.0} & 92.0 & 100.0 & 100.0 & 100.0 \\
        Hallucination Rate (\%) $\downarrow$ & 55.0 & 26.0 & \textbf{5.0} & 47.0 & 12.0 & \textbf{11.0} \\
        \bottomrule
    \end{tabular}
\end{table}
%===============================================

% \begin{figure}[t]
% \centering
% \begin{tikzpicture}
% \begin{groupplot}[
%     group style={group size=2 by 1, horizontal sep=2cm},
%     ymin=0, ymax=100,
%     width=6.5cm,
%     height=4.8cm,
%     xtick={1,2,3},
%     xticklabels={B1,B2,B3},
%     ytick={0,20,40,60,80,100},
%     ymajorgrids,
%     grid style={dashed, gray!40},
%     tick label style={font=\footnotesize},
%     ylabel style={font=\small},
%     title style={font=\small\bfseries},
%     legend style={font=\footnotesize, at={(0.5,-0.25)}, anchor=north, legend columns=2}
% ]

% \nextgroupplot[
%     ybar,
%     bar width=15pt,
%     title={Case I},
%     ylabel={Percentage (\%)}
% ]
% \addplot[fill=gray!35, draw=black!70] coordinates {(1,45) (2,73) (3,85)};
% \addplot[fill=black!75, draw=black] coordinates {(1,55) (2,26) (3,5)};

% \nextgroupplot[
%     ybar,
%     bar width=15pt,
%     title={Case II}
% ]
% \addplot[fill=gray!35, draw=black!70] coordinates {(1,23) (2,76) (3,81)};
% \addplot[fill=black!75, draw=black]   coordinates {(1,47) (2,12) (3,11)};
% \legend{Answer Correctness (\%), Hallucination Rate (\%)}

% \end{groupplot}
% \end{tikzpicture}

% \caption{Answer correctness and hallucination rate across baselines and case studies ($n = 100$ queries each).}
% \label{fig:performance_bars}
% \end{figure}

%===============================================

The directional pattern B1\,$<$\,B2\,$<$\,B3 is preserved for answer correctness and hallucination reduction across both domains (Table~\ref{tab:performance_combined}). Detailed error attribution is provided in Section~\ref{subsec:error_analysis}.
 
Statistical significance was assessed using McNemar's test on matched paired judgements ($n = 100$). Results are summarised in Table~\ref{tab:mcnemars_tests}. In both case studies, the B1$\rightarrow$B2 improvement in answer correctness is highly significant ($p < 0.001$). The additional B2$\rightarrow$B3 correctness gain is significant in Case~I ($\chi^2 = 4.65$, $p = 0.031$) but does not reach significance in Case~II ($\chi^2 = 0.84$, $p = 0.359$). Hallucination reductions from B1$\rightarrow$B3 are highly significant in both cases ($p < 0.001$). The B2$\rightarrow$B3 hallucination change in Case~II (12\%$\rightarrow$11\%) amounts to a single query and was not tested separately.


% SIMPLIFIED TABLE WITHOUT ns
% \begin{table}[t]
% \centering
% \small
% \caption{Statistical significance of baseline improvements (McNemar's test, $n=100$ paired queries). $\Delta$ = percentage-point change; 95\% CI computed via Wilson method on matched proportions.}
% \label{tab:mcnemars_tests}
% \setlength{\tabcolsep}{4pt}
% \renewcommand{\arraystretch}{1.1}
% \begin{tabular}{@{}llrrrr@{}}
% \toprule
% \textbf{Case} & \textbf{Comparison} & \textbf{$\Delta$ (pp)} & \textbf{$\chi^2$} & \textbf{$p$} & \textbf{95\% CI} \\
% \midrule
% \multirow{4}{*}{Case~I}
%   & B1$\to$B2 Correctness   & $+28$ & 15.85 & $<0.001$ & [16,\;40] \\
%   & B2$\to$B3 Correctness   & $+12$ &  4.65 & 0.031    & [2,\;22]  \\
%   & B1$\to$B3 Correctness   & $+40$ & 34.57 & $<0.001$ & [30,\;50] \\
%   & B2$\to$B3 Hallucination & $-21$ & 13.79 & $<0.001$ & [$-31$,\;$-11$] \\
% \midrule
% \multirow{4}{*}{Case~II}
%   & B1$\to$B2 Correctness   & $+53$ & 40.36 & $<0.001$ & [41,\;65] \\
%   & B2$\to$B3 Correctness   &  $+5$ &  0.84 & 0.359    & [$-3$,\;13] \\
%   & B1$\to$B3 Correctness   & $+58$ & 45.12 & $<0.001$ & [46,\;70] \\
%   & B1$\to$B3 Hallucination & $-36$ & 25.52 & $<0.001$ & [$-48$,\;$-24$] \\
% \bottomrule
% \end{tabular}
% \end{table}





% FULL TABLE
\begin{table}[t]
\centering
\small
\caption{Statistical significance of baseline improvements (McNemar's test, $n=100$ paired queries). $\Delta$ = percentage-point change; 95\% CI on matched proportions (Wilson method). Rows marked \textit{ns} do not reach $p<0.05$.}
\label{tab:mcnemars_tests}
\setlength{\tabcolsep}{4pt}
\renewcommand{\arraystretch}{1.1}
\begin{tabular}{@{}llrrrr@{}}
\toprule
\textbf{Case} & \textbf{Comparison} & \textbf{$\Delta$ (pp)} & \textbf{$\chi^2$} & \textbf{$p$} & \textbf{95\% CI} \\
\midrule
\multirow{6}{*}{Case~I}
  & B1$\to$B2 Correctness   & $+28$ & 15.85 & $<0.001$ & [16,\;40]       \\
  & B2$\to$B3 Correctness   & $+12$ &  4.65 & 0.031    & [2,\;22]        \\
  & B1$\to$B3 Correctness   & $+40$ & 34.57 & $<0.001$ & [30,\;50]       \\
  & B1$\to$B2 Hallucination & $-29$ & 16.68 & $<0.001$ & [$-41$,\;$-17$] \\
  & B2$\to$B3 Hallucination & $-21$ & 13.79 & $<0.001$ & [$-31$,\;$-11$] \\
  & B1$\to$B3 Hallucination & $-50$ & 48.02 & $<0.001$ & [$-60$,\;$-40$] \\
\midrule
\multirow{6}{*}{Case~II}
  & B1$\to$B2 Correctness   & $+53$ & 40.36 & $<0.001$ & [41,\;65]       \\
  & B2$\to$B3 Correctness   &  $+5$ &  0.84 & \textit{ns}~0.359 & [$-3$,\;13] \\
  & B1$\to$B3 Correctness   & $+58$ & 45.12 & $<0.001$ & [46,\;70]       \\
  & B1$\to$B2 Hallucination & $-35$ & 22.67 & $<0.001$ & [$-47$,\;$-23$] \\
  & B2$\to$B3 Hallucination &  $-1$ &  0.00 & \textit{ns}~1.000 & [$-7$,\;5]  \\
  & B1$\to$B3 Hallucination & $-36$ & 25.52 & $<0.001$ & [$-48$,\;$-24$] \\
\bottomrule
\end{tabular}
\end{table}

%============================================

%===============================================
\subsubsection{Operational Cost}
\label{subsubsec:operational_cost}
%============================================

\begin{table}[h]
    \centering
    \caption{Operational cost comparison across both case studies ($n = 100$ queries each).}
    \label{tab:operational_cost_combined}
    \small
    \setlength{\tabcolsep}{5pt}
    \renewcommand{\arraystretch}{1.1}
    \begin{tabular}{@{}lrrrrrr@{}}
        \toprule
        & \multicolumn{3}{c}{\textbf{Case~I}} & \multicolumn{3}{c}{\textbf{Case~II}} \\
        \cmidrule(lr){2-4}\cmidrule(lr){5-7}
        \textbf{Metric} & \textbf{B1} & \textbf{B2} & \textbf{B3} & \textbf{B1} & \textbf{B2} & \textbf{B3} \\
        \midrule
        Avg.\ Latency (s)    & 12.24 & 0.48  & 13.65 & 5.25  & 2.55  & 12.84 \\
        Avg.\ Function Calls & 0     & 1     & 5.2   & 0     & 1     & 4.2   \\
        \bottomrule
    \end{tabular}
\end{table}

Table~\ref{tab:operational_cost_combined} reports per-query latency and function call counts. The RCP framework (B3) incurs the highest per-query latency in both cases due to its multi-stage control loop with multiple LLM reasoning steps and iterative tool execution. SQL-backed retrieval (B2) achieves the lowest latency (0.48\,s in Case~I and 2.55\,s in Case~II) because it executes a single deterministic SQL join.

In Case~I, B1 latency (12.24\,s) is comparable to B3 (13.65\,s) because B1's embedding cost over the full catalogue nearly matches B3's multi-step overhead. In Case~II, the gap is larger (B1: 5.25\,s vs.\ B3: 12.84\,s) because B1 operates over a smaller corpus requiring less embedding computation. The latency--reliability trade-off is inherent to the RCP design: validation gates and persistent state tracking are the mechanisms that nearly eliminate hallucination, and each requires additional computation. In domains where correctness is a compliance or safety requirement, the additional latency of B3 is the cost of auditability.
 
%============================================
\subsection{Error Analysis}
\label{subsec:error_analysis}
%============================================
Appendix Tables~\ref{tab:qs_hydroscand} and~\ref{tab:qs_saab} list sample queries from the annotated sets used in Section~\ref{sec:evaluation}. The test queries, representative of practical use scenarios, were authored by domain engineers at the respective companies who had no involvement in framework development, mitigating data leakage between query authoring and system design.

Each query is assigned a difficulty tier reflecting the number of reasoning steps and source documents required for a correct answer (\emph{Easy}: single table lookup; \emph{Medium}: cross-field inference or semantic matching; \emph{Hard}: multi-document reasoning or domain-specific disambiguation). Tier distributions are reported in the appendix tables; a stratified performance breakdown by tier is reserved for future work.

The \emph{Ground truth} column records the correct answer as determined by the domain expert with reference to the source documents. For Case~I, ground-truth answers were verified against the Hydroscand Produktbok PDF (hydraulic hose and press coupling chapters), whereas ground-truth answers for Case~II were verified against the applicable revision of the SAAB AB technical source document.

The \emph{Error class} column identifies the primary failure mode of the baselines on each query, drawn from the following taxonomy:
\begin{enumerate}[label=\roman*.]
    \item Hallucination: the model asserts a specific value (article number, torque, revision) not present in any retrieved chunk.
    \item Partial answer: a correct product family or category is identified but the required specificity (dimension, standard designation, per-revision caveat) is absent.
    \item Field-name mismatch: the correct source table is retrieved but the query term does not match the column header literally, causing the retrieval step to return no evidence.
    \item Parser error: the source document structure could not be reliably parsed, resulting in malformed or missing field values.
\end{enumerate}

The B1$\rightarrow$B2 transition reveals that the dominant failure mode in Naive RAG is retrieval fragmentation: chunks rarely contain complete specification records, leading the LLM to hallucinate values for missing fields (hallucination rate 55\% in Case~I, 47\% in Case~II). Structured SQL retrieval reduces hallucination to 26\% in Case~I and 12\% in Case~II by delivering complete product records with preserved field semantics, lifting answer correctness from 45\% to 73\% (Case~I) and from 23\% to 76\% (Case~II). Residual B2 errors fall into two classes: keyword-search failures on semantic queries without explicit product codes (where LIKE-pattern fallback retrieves irrelevant records), and generation errors when multiple products are returned and the LLM conflates their attributes.

The B2$\rightarrow$B3 transition shows further improvement. Answer correctness rises to 85\% (Case~I) and 81\% (Case~II) while hallucination is nearly eliminated (5\% in Case~I, 11\% in Case~II). The RCP's multi-step planning enables it to handle both direct lookups and comparison queries through specialised strategies (\emph{Specification Lookup}, \emph{Multi-Product Comparison}), while validation gates prevent the LLM from fabricating plausible but unsupported values. The residual hallucination in B3 appears to originate primarily in the LLM synthesis step, where validated attributes are converted to natural language. This verify-then-summarise behaviour is the intended design outcome: in engineering contexts, an explicit qualified answer is preferable to a confident but incorrect output.

\subsubsection{Metric-Level Anomalies}
\label{subsubsection:metric_anomalies}
%============================================
Two metrics show non-monotone behaviour that warrants explanation.
 
\textbf{Citation accuracy:} The RCP's validation gates enforce schema completeness, attribute presence, and unit correctness, but do not include an explicit \emph{citation-completeness} check. Consequently, the RCP may produce correct answers grounded in verified attributes without citing the specific catalogue page or table cell, while Naive RAG inherently propagates chunk metadata (document name, page, offset). Adding a citation-provenance predicate to strategy-level sufficiency conditions is a planned extension.
 
\textbf{Unit fidelity:} The drop from B2 (93\%) to B3 (92\%) in Case~I amounts to a single query on $n = 100$ and is within sampling error. The directional explanation is nonetheless instructive. B2 retrieves a complete product record via SQL and passes it directly to the LLM, which echoes whatever unit representation is present in the source data---including cases where the source already uses a canonical format. B3, by contrast, routes all values through \texttt{func\_normalize\_units} and enforces a unit-validation predicate before synthesis. When normalisation produces an ambiguous or unresolvable result (e.g., a non-standard thread notation or a missing unit field), the validation gate returns a qualified non-answer rather than propagating the raw value. This reduces hallucination at the cost of a marginally lower raw fidelity score: some unit values that B2 passes through unchecked would have been accepted as correct under the evaluation rubric, whereas B3 explicitly rejects them. Unit fidelity is perfect across all baselines in Case~II, reflecting the more standardised unit notation in aerospace documentation. Adding tolerance bands or graduated unit-check predicates is a planned extension.

% DISCUSSION AND LIMITATIONS
%============================================
\section{Discussion and Limitations}
\label{sec:Discussion}
%============================================
The evaluation demonstrates a clear trade-off: the RCP improves reliability and traceability of LLM-based retrieval in engineering settings, but at measurable computational cost. Compared to the SQL-backed baseline~(B2), the multi-stage control loop adds approximately 13\,s per query in Case~I and 10\,s in Case~II (Table~\ref{tab:operational_cost_combined}), reflecting additional function calls, validation checkpoints, and persistence of intermediate artefacts.
 
In both case studies, the baselines occasionally produced confident but incorrect outputs, particularly in revision-controlled scenarios. The RCP's verify-then-summarise behaviour is the mechanism that accounts for the B2-to-B3 correctness increment and the near-elimination of hallucination: explicit failure when sufficiency conditions are unmet is preferable to silent misinformation in contexts where unit or factual errors propagate downstream. This confirms the motivation advanced by Leng et al.~\cite{leng2026}, who argue that agentic systems require explicit validation mechanisms rather than black-box reasoning, and extends the self-refinement approach of Madaan et al.~\cite{madaan2023selfrefine} from free-form language feedback to typed, schema-level predicates over structured artefacts.
 
The performance gap between cases reflects the greater complexity of Case~II, which requires cross-document reasoning with revision scoping and multimodal evidence fusion, increasing the surface area for residual LLM synthesis errors that pass validation gates. Closing this gap would likely require richer strategy-level predicates (e.g., cross-source consistency checks) or stronger base models.

In domains where outputs inform compliance or design decisions, correctness and auditability are primary requirements, and the observed latency overhead may be acceptable. A practical consequence of the RCP's relational design is that inspection, error attribution, and replay reduce to database queries rather than bespoke logging, lowering the engineering effort associated with debugging, evaluation, and governance in production deployments.
 
Because validation predicates operate over structured artefacts independently of the LLM (they do not depend on any model-specific output format or internal representation), the architectural benefits are expected to hold across model families. However, the current evaluation uses a single compact model (Llama3.2:8b) for reasoning, and the absolute correctness and hallucination rates will vary with model capability. Cross-model evaluation remains future work.

%============================================
\vspace{1.0em}
\noindent \textbf{Limitations.}
%============================================
 
\begin{description}[style=unboxed, leftmargin=0pt, labelindent=0pt, font=\itshape]
 
\item[Upstream extraction quality.]
System reliability depends on document parsing. OCR errors, poor table reconstruction, and ambiguous typography propagate into the harvested database and constrain downstream reasoning. The RCP prevents unsupported synthesis but cannot correct malformed source data.
 
\item[Predicate tuning.]
Validation predicates require careful calibration. Overly strict predicates increase retries and latency; overly permissive predicates allow subtle errors to propagate undetected. Predicate design thus represents an engineering trade-off between computational cost and error tolerance.
 
\item[Domain scope.]
The evaluation was limited to two industrial domains. Although the title characterises the architecture as domain-agnostic, this claim is supported by architectural argument: the separation of orchestration logic from domain-specific strategy and function libraries, rather than by empirical evidence across a broad range of document genres. Broader benchmarking across legal contracts, medical device documentation, or legacy scanned manuals would strengthen the generalisability claim.
 
\item[Extensibility overhead.]
The current implementation models agent capabilities as canonical functions registered in a structured function library. Although this supports deterministic execution, it introduces integration overhead when extending to new domains. Emerging skill-based capability abstractions~\cite{agentskills2024} that define modular, discoverable capability manifests could reduce this friction while preserving governance guarantees.
 
\item[Cross-session learning.]
The framework does not adapt strategy selection or function parameters based on historical execution outcomes; each session starts from a clean state. Incorporating cross-session learning, for example, biasing strategy selection toward templates that have historically succeeded for structurally similar queries, is planned future work.
 
\item[Annotation protocol.]
Ground-truth answers were determined by the domain engineers per case study. No formal inter-annotator agreement was assessed. While the technical specificity of engineering queries limits ambiguity, independent second-annotator verification would strengthen the evaluation.
 
\end{description}


% CONCLUSION
%============================================
\section{Conclusions and Future Work}
\label{sec:conclusions_and_future_work}
%============================================
This study introduced the Relational Control Plane (RCP), a SQL-backed agentic architecture that persists orchestration state as queryable relational records. The architecture enforces a six-stage verify-then-summarise control loop over engineering document retrieval, in which synthesis is permitted only after retrieved evidence satisfies explicit validation constraints.
 
Evaluation across 200 annotated queries spanning two industrial domains---hydraulic product catalogues and revision-controlled aerospace documentation---showed statistically significant improvements. Notably, these gains were achieved with a compact 8B-parameter model, suggesting the architecture rather than model scale drives performance. Answer correctness increased from 45\% to 85\% in Case~I and from 23\% to 81\% in Case~II, while hallucination decreased from 55\% to 5\% and from 47\% to 11\% in the respective domains. The framework preserved engineering-critical fidelity in unit correctness, applicability scope, and provenance granularity, while transforming potential hallucinations into explicit, auditable failures.
 
Future work will focus on (i)~expanding validation toward cross-source consistency checking, including conflict detection and supersession-aware reconciliation across revisions; (ii)~improving multimodal grounding for tables and figures through tighter cell- and region-level anchoring; (iii)~incorporating cross-session learning to adapt strategy selection based on historical execution outcomes; (iv)~developing evaluation suites that score provenance completeness, revision fidelity, and replay determinism alongside answer correctness; and (v)~cross-model evaluation to assess whether architectural gains generalise across model scales and commercial-grade systems.
 
The results indicate that treating orchestration as structured, persistent data provides a practical and extensible foundation for deploying reliable generative AI systems in engineering knowledge management.
%==============================================================================


%==============================================================================
% CLOSING
%==============================================================================
\section*{Acknowledgements}
The authors gratefully thank the industrial partners who provided representative technical documentation and valuable domain feedback.

\noindent During the preparation of this manuscript, the authors used ChatGPT (GPT-[5.3 ], OpenAI, accessed February 2026) for code refinement and language editing to improve clarity and readability. All outputs were reviewed, validated, and edited by the authors to ensure accuracy and compliance with scientific standards.
%_______________________________________________________________________________
\section*{Author Contributions}
CRediT: \textbf{Oscar Ikechukwu:} Formal analysis, Methodology, Software, Validation, Visualization, Writing -- original draft, Writing -- review \& editing; \textbf{Mehdi Tarkian:} Conceptualization, Methodology (framework and architecture design), Formal analysis (theoretical framework development), Supervision, Writing -- review \& editing; \textbf{Sanjay Nambiar:} Investigation, Writing -- review \& editing; \textbf{Marie Jonsson:} Writing -- review \& editing; \textbf{Christoffer Brax:} Case study II -- Data curation.

%_______________________________________________________________________________
\section*{Disclosure Statement}
No potential conflict of interest was reported by the author(s).
%_______________________________________________________________________________
\section*{Funding}
\noindent This work was supported by Vinnova; Sweden's Innovation Agency under Grant [2024-01420] - Data Automation and Retrieval Technology (DART) as part of the research program 'Advanced digitalization - Enabling technologies'.
%_______________________________________________________________________________
\section*{ORCiD}
\textit{Oscar Ikechukwu:} \orcidlink{0009-0000-4905-2344}  https://orcid.org/0009-0000-4905-2344\\
\textit{Mehdi Tarkian:} \orcidlink{0009-0003-6305-8621}  https://orcid.org/0009-0003-6305-8621\\
\textit{Sanjay Nambiar:} \orcidlink{0000-0003-1745-3869}  https://orcid.org/0000-0003-1745-3869\\
\textit{Marie Jonsson:} \orcidlink{0000-0002-6079-2359} https://orcid.org/0000-0002-6079-2359\\
\textit{Christoffer Brax:} \orcidlink{0000-0002-2161-164X} https://orcid.org/0000-0002-2161-164X
%_______________________________________________________________________________

\section*{Data Availability Statement}
 The code for the RCP-Framework is openly available on GitHub at \url{https://github.com/oscik559/RCP-Framework}.
The data supporting Case~I (Hydroscand hydraulic catalogue) available in the same repo. However, the data supporting Case~II are subject to third-party restrictions and were used under confidentiality agreement with SAAB AB; they are not 
publicly available. Requests may be directed to the SAAB AB correspondent listed in the author notes.

% The implementation code supporting both case studies is available on GitHub at \url{https://github.com/oscik559/RCP-Framework}. This includes the full extraction pipeline, the RCP agentic framework, evaluation baselines (B1--B3), annotated query sets, and result files for Case~I and Case~II. The product database for Case~I (Hydroscand hydraulic catalogue) is included in the repository. The raw source data supporting Case~II are subject to third-party restrictions and were used under a confidentiality agreement with SAAB AB; they are not publicly available. Requests may be directed to the SAAB AB correspondent listed in the author notes.
%==============================================================================
% References
%==============================================================================
\bibliographystyle{tfnlm}
\bibliography{references}
%==============================================================================
% APPENDICES
%==============================================================================
\clearpage
% \section*{Appendices}
\appendix
% Configure appendix numbering (ADD THIS!)
\renewcommand{\thesection}{\Alph{section}}
\renewcommand{\thesubsection}{\thesection.\arabic{subsection}}
\renewcommand{\thetable}{\thesection.\arabic{table}}
\setcounter{section}{0}

% Now input your appendix files
\newpage
\input{appendix/annotated_query_sets}           % Will be Appendix C
\input{appendix/prompts}                % Will be Appendix A  
% \input{appendix/sample_query_sets}              % Will be Appendix B
% \newpage
\input{appendix/detailed_execution_traces}                % Will be Appendix D
%_______________________________________________________________________________

\end{document}